% !TeX program = tectonic
\documentclass[11pt,a4paper]{article}
\usepackage{fontspec}
\usepackage[english]{babel}
\babelfont{rm}[Language=Default]{Liberation Serif}
\babelfont[Language=Default]{sf}{Carlito}
\babelfont[Language=Default]{tt}{DejaVu Sans Mono}
\usepackage{amsmath,amssymb,amsthm}
\usepackage{geometry}
\geometry{a4paper, top=2.4cm, bottom=2.4cm, left=2.4cm, right=2.4cm}
\usepackage{booktabs,tabularx,enumitem,xcolor,tcolorbox}
\tcbuselibrary{breakable,skins}
\definecolor{linkblue}{HTML}{1F4E79}
\definecolor{statgreen}{HTML}{2F6B3A}
\definecolor{goldacc}{HTML}{8B7E5A}
\usepackage[colorlinks=true,linkcolor=linkblue,urlcolor=linkblue,unicode=true]{hyperref}
\theoremstyle{plain}
\newtheorem{theorem}{Theorem}
\newtheorem{lemma}[theorem]{Lemma}
\newtheorem{proposition}[theorem]{Proposition}
\newtheorem{corollary}[theorem]{Corollary}
\theoremstyle{definition}
\newtheorem{definition}[theorem]{Definition}
\theoremstyle{remark}
\newtheorem{remark}[theorem]{Remark}
\newtcolorbox{statusbox}[1][]{colback=statgreen!5,colframe=statgreen!75,
  title={\textbf{Summary of results:} #1},fonttitle=\small,boxrule=0.7pt,arc=2pt,breakable}
\newtcolorbox{databox}[1][]{colback=goldacc!6,colframe=goldacc!80,
  title={\textbf{Protocol data:} #1},fonttitle=\small,boxrule=0.7pt,arc=2pt,breakable}
\newtcolorbox{verifybox}[1][]{colback=linkblue!5,colframe=linkblue!80,
  title={\textbf{Verification:} #1},fonttitle=\small,boxrule=0.7pt,arc=2pt,breakable}
\widowpenalty=10000
\clubpenalty=10000
\setlength{\parskip}{0.35em}
\newcommand{\CC}{\mathbb{C}}
\newcommand{\RR}{\mathbb{R}}
\newcommand{\ZZ}{\mathbb{Z}}
\newcommand{\QQ}{\mathbb{Q}}
\newcommand{\Ga}{\Gamma}
\newcommand{\Bfun}{\mathrm{B}}
\newcommand{\im}{\operatorname{Im}}
\newcommand{\re}{\operatorname{Re}}
\newcommand{\lcm}{\operatorname{lcm}}
\newcommand{\SNF}{\operatorname{SNF}}
\newcommand{\Jac}{\operatorname{Jac}}
\newcommand{\rank}{\operatorname{rank}}
\newcommand{\Ind}{\operatorname{Index}}
\newcommand{\disc}{\operatorname{disc}}
\begin{document}
\begin{center}
{\LARGE\bfseries Certificate B: Closing the Kernel 29}\\[0.4em]
{\large exact comparison with periods and the Hodge test}\\[0.6em]
{\normalsize Isaev Iskhak Khamzatovich}\\[0.2em]
{\normalsize The <<Dynamic Principle>> Program --- the \texttt{hodge-laboratory}}\\[0.15em]
{\small Standalone theorem monograph · Монография 10/16}
\end{center}
\vspace{0.4em}
{\color{linkblue}\hrule height 0.8pt}
\vspace{0.8em}

\section{Setting}

The cycle certifier (certificate A) presents a ``real figure'' ---
an explicit divisor with equations. Its blind spot --- the
dimension-$29$ pairing kernel --- is closed by certificate~B with
periods. Together A and B guarantee: the divisor class is exactly
determined on the whole cohomology.

\begin{theorem}[B1 --- structure of the kernel]\label{th:B1}
Let $G$ be the exact pairing matrix of $49$ measurements with the
rank-$20$ algebraic sublattice. Then $\dim\ker G=29$, the kernel
consists exactly of the cohomological relations, and certificate A
is complete on the N\'eron--Severi sublattice. The blind spot
$29+2\to0$: the two transcendental directions are checked by
periods to $<10^{-90}$.
\end{theorem}

\begin{proof}
The rank of $G$ is $20$ (the K3 invariants), the measurement space
has dimension $49$, hence $\dim\ker G=29$. The kernel is the
orthogonal complement of the algebraic sublattice; by the Hodge
index theorem (Shioda, Fermat quartic) it consists exactly of the
cohomological relations: the classes of $H^{2,0}\oplus H^{0,2}$ and
the anti-algebraic part of $H^{1,1}$. The period Hodge test: the two
$(2,0)$ directions are checked by integrals; values $<10^{-90}$ are
identified with zero by the exact analytics.
\end{proof}

\begin{databox}[certificate B numbers]
The divisor $Z=\tfrac32 L_1(0,0)-\tfrac12 L_2(1,1)+\tfrac54 h$ ---
all pairings exact over $\QQ$. Torus: target class
$5\cdot[\mathrm{pt}]$; degree certificate exact; period lattice
$\ZZ+i\ZZ$ exact; Abel--Jacobi sum $5\pi^2/32+15/4$; verdict:
accepted fully. K3: cup-form signature $(2,19)$; $51$ measurements
certify the whole $22$-dimensional $H^2$; assembly
$[Z]=[target]$ in $H^2(X,\QQ)$; verdict: accepted.
\end{databox}

\begin{remark}[the closing logic]
$29$ relations are cohomological (structurally described by the
Hodge index theorem), $2$ dimensions are transcendental (closed by
periods). The blind spot $29+2$ is reduced to zero: every
cohomology direction is either measured or structurally described.
The coincidence of the factor $29$ with $406=2\cdot7\cdot29$ is an
independent observation, recorded as such.
\end{remark}

\begin{statusbox}[summary]
Certificate B accepted: the kernel $29$ closed by Theorem B1, the
transcendental directions by periods; the class assembly is exact in
$H^2(X,\QQ)$.
\end{statusbox}

\begin{verifybox}[reproduction]
\texttt{python3 laboratory.py} $\to$ ``Certificates $\to$ B'': the
assembly and the Hodge test --- seconds; the JSON report in
\texttt{reports/}.
\end{verifybox}

\vfill
{\color{linkblue}\hrule height 0.6pt}
\vspace{0.5em}
{\small\color{gray}
License: individual exclusive license of Isaev Iskhak Khamzatovich.
All rights to the computations, formulas, and texts belong to the author.
Verification: \texttt{python3 laboratory.py} (``Stands''/``Certificates''),
Lean: \texttt{verification/lean/HodgeLaboratory.lean}.
}
\end{document}
